%%%%%%%%%%%%%%%%%%%%%%%%%%%%%%%%%
%
%       EXERCÍCIO
%
%%%%%%%%%%%%%%%%%%%%%%%%%%%%%%%%%

\definevalorquestao

\begin{exercicioBanco}[\valorquestao]
O número de acidentes de trabalho por ano registrados por funcionários de uma indústria foi modelado pela variável aleatória $X$, cuja função de probabilidade é dada por:
\[
\begin{array}{c|ccc}
x & 0 & 1 & 2 \\ \hline
P(X = x) & 0{,}6 & 0{,}3 & 0{,}1
\end{array}
\]

\noindent
Para estimar a média de acidentes $\mu = E(X)$, sorteou-se uma amostra aleatória com reposição de tamanho $n = 2$, representada por $(X_1, X_2)$, e propôs-se o seguinte estimador:
\[
\hat{\mu} = \frac{X_1 + 2X_2}{3}
\]

\begin{enumerate}[(a)]
    \item Calcule o valor verdadeiro da média populacional $\mu = E(X)$.
    \item Obtenha a distribuição de probabilidade do estimador $\hat{\mu}$.
    \item Verifique se o estimador $\hat{\mu}$ é não viciado para $\mu$.
\end{enumerate}
\end{exercicioBanco}

